% TREŚĆ ROZDZIAŁU: zastąp wskazówki opisem własnej koncepcji i projektu.
% Rysunki przechowuj w katalogu figures/ rozdziału, podpisuj przez \caption
% i oznaczaj przez \label, aby później przywołać je poleceniem \cref.
\chapter{Koncepcja i projekt rozwiązania}
\label{chap:projekt}

\section{Analiza wariantów}
Przedstaw rozważane rozwiązania, kryteria porównania i uzasadnienie wyboru. Nie ograniczaj uzasadnienia do stwierdzenia, że wariant był najłatwiejszy.

\section{Architektura rozwiązania}
Pokaż moduły, granicę systemu, przepływ informacji lub energii oraz interfejsy. Każdy schemat omów w tekście.

\section{Projekt szczegółowy}
Opisz algorytmy, modele, struktury danych, interfejsy, podzespoły, schematy lub parametry niezbędne do wykonania rozwiązania. Liczby z jednostkami zapisuj np. jako \qty{230}{\volt}.

\section{Plan weryfikacji}
Powiąż przypadki testowe z wymaganiami. Dla każdego testu określ dane wejściowe, procedurę, oczekiwany rezultat i warunek zaliczenia.
